%% ==========================================================================
%%  9  Conclusion and Future Work
%%
%%  Mirrors R3's Section 5.  The two "don't go here" results are the reason
%%  this section is worth reading before starting work.
%% ==========================================================================
\section{Conclusion and Future Work}
\label{sec:conclusion}

We have not settled Conjecture~\ref{conj:main} with a proof, but
\Cref{sec:approach6} now carries a construct-and-repair algorithm with
\textbf{zero observed failures across $389{,}215$ test instances}, including
every hard instance in this project. That is the headline result: strong
constructive evidence for the conjecture, from an algorithm precise enough that
the remaining gap to a proof is a single, sharply stated combinatorial claim
(Conjecture~\ref{conj:algorithm-succeeds}).

Getting there required ruling out three strategies a reader would reasonably try
first: item-by-item insertion, the strategy that works in~\cite{BKNS22}
(Theorem~\ref{thm:obstruction}, and again in the pure goods coordinates of
\Cref{sec:insertion-fails-goods}); restriction to utilitarian-optimal
allocations, together with \emph{every} refinement of that set
(Proposition~\ref{prop:msw-false}, Corollary~\ref{cor:no-refinement}); and
encoding the coverage constraint into the numbers, at either the valuation or
the algorithm level (\Cref{sec:approach4}). Alongside those sit a small set of
positive results --- binary additive costs (Theorem~\ref{thm:binadd}, now
subsumed by~\cite{LMS26}), identical costs (Theorem~\ref{thm:identical}), two
agents --- and structural tools that say precisely how the chore problem differs
from its goods mirror.

\paragraph{The lesson that ties the whole investigation together.}
Every insertion-based strategy failed for the same reason, in three different
coordinate systems: chores directly (\Cref{sec:approach1}), the replica/peel
frame (\Cref{sec:approach3}), and even the pure goods frame
(\Cref{sec:insertion-fails-goods}). Building an allocation one item at a time
fixes decisions before their consequences are visible, and by the time a bad
grouping shows itself, insertion has no way to undo it --- only permutation
between agents, never a genuine re-split of a bundle's contents, is available.
The algorithm that finally worked abandoned that template entirely:
\emph{construct} a whole candidate partition directly, then \emph{repair} it by
moves that can freely undo an earlier bad grouping. That is the one idea from
this project worth carrying into any future attempt on a similar problem.

The three negative results are more useful together than separately, because
they fail in the same way. Insertion fixes an owner before the consequences are
visible; a selection rule fixes a criterion before knowing which allocation it
will pick out; a reweighting fixes a penalty before knowing who will hold the
duplicate. Each is \emph{a rule fixed ex ante against a quantity determined ex
post} (\Cref{sec:holderblind}). What survives is procedures that decide late ---
which is exactly the design of the peel frame, where ownership is settled only
by the last move touching a chore.

Both negative results are failures of \emph{timing}, and the constraint they
jointly impose is that \emph{a correct algorithm must be free to move
already-allocated chores between bundles, and free to give up total cost while
doing so}. The replica transform of \Cref{sec:approach3} is the response:
Theorem~\ref{thm:replica} turns the instance into a goods instance with an
identical envy graph, Corollary~\ref{cor:coverage} reduces the conjecture to the
single constraint that no agent be relieved of a chore twice, and reading that
dynamically defers every ownership decision to the last move touching a chore,
so there is nothing to re-open. This is currently the most promising route.
It is not known to work --- Proposition~\ref{prop:deadends} shows the state
space has genuine dead ends, and Remark~\ref{rem:converse} notes that the
reduction is sufficient but not known to be complete --- and what is missing is
a peel rule that provably never enters a deadlock.

Every construction here that succeeds does so by exhibiting a per-agent
potential $\phi$ with $\edgew{\alloc}(i,j) \le \phi_i - \phi_j$, which
telescopes along paths. Finding such a potential for general negative
dichotomous costs, or proving none exists, remains the problem; the balance
signal of Remark~\ref{rem:balance} is the current best guess at its shape.

\paragraph{Directions.}
\begin{enumerate}
  \item \textbf{Prove Conjecture~\ref{conj:algorithm-succeeds}}
    (\Cref{sec:construct-repair-remains}). This is now the top-ranked direction
    by a wide margin: it is the only open statement in the project backed by a
    zero-failure sweep of $389{,}215$ instances, and it is purely combinatorial
    --- cardinality-balanced partitions and single-item swaps, no chores, no
    replica, no schedule. The $14$ known local optima are the concrete obstacle
    a proof must explain: they show bad local optima exist in the swap
    neighbourhood, but every one found so far was escaped by restarting from a
    different starting partition. A proof needs either a potential that a
    restart-reachable partition always decreases, or a precise characterisation
    of the bad optima sharp enough to name a starting rule that avoids them.
  \item \textbf{The type-ordered peel search} (\Cref{sec:typeorder},
    \Cref{sec:peelsearch}). Conjecture~\ref{conj:h1pp} has the best evidence of
    anything currently open and a much smaller state space than
    Conjecture~\ref{conj:h1prime}. Run the fixed-point computation over the full
    $n = m = 3$ dichotomous family; a bad root anywhere kills the approach of
    \Cref{sec:approach3}, a clean sweep is strong evidence for it. Cheapest
    decisive experiment available. Note what a recipient rule must now satisfy:
    by Theorem~\ref{thm:f2} the recipient in the hard case has zero marginal, and
    by Proposition~\ref{prop:f4} $\argmax_i \pathw{}(i)$ is insufficient even
    under the type-ordering, so the rule has to read the residual workload.
  \item \textbf{Choose the paid set first} (Proposition~\ref{prop:twotier},
    Proposition~\ref{prop:twotier-equiv}). The target is exactly: an
    envy-freeable allocation together with a bipartition of the agents such that
    all strict envy runs from one side to the other. This is a condition on a
    bipartition rather than on a schedule, so it inherits none of the dead ends
    of \Cref{sec:approach3} and none of the ex-ante failures of
    \Cref{sec:approach4}. It is the one direction that four separate lines of
    attack have now pointed at and nobody has tried.
  \item \textbf{Agent induction with a path-based invariant}
    (\Cref{sec:approach5}). The arc-based invariant is refuted, so any revival
    must carry path weights, which are not assignment-invariant. The equality
    graph of~\cite[\S4.3.1]{LMS26} is the natural candidate for what to carry.
    Harder than it looked a session ago, but not closed.
  \item \textbf{Characterise the dead ends} (Remark~\ref{rem:balance}). Test
    across that family whether every dead end is a state already committed to an
    unbalanced terminal. If so, the balance rule is the missing ingredient and
    can be stated precisely.
  \item \textbf{Settle the converse} (Remark~\ref{rem:converse}). Is a good
    terminal state, when one exists, always reachable by a legal schedule? A
    positive answer upgrades the peel frame from a sufficient reformulation to a
    complete one, and makes the search of item~1 a direct test of
    Conjecture~\ref{conj:main}.
  \item \textbf{Adversarial instance generation} (\Cref{sec:samplerbias}). The
    randomised evidence is drawn from a distribution that essentially never
    produces the structured extremal functions from which both negative results
    are built. This is the cheapest way to find out whether
    Conjecture~\ref{conj:main} is true at all.
  \item \textbf{Exchange paths rather than single transfers}
    (\Cref{sec:tiebreaks}). Single-chore local search has bad local minima. The
    move set should match the exchange structure of dichotomous costs, as
    Yankee Swap and matroid union do in the matroid-rank
    literature~\cite{BEF21}; and by Remark~\ref{rem:msw-deadend} the search must
    range over all allocations, not the utilitarian-optimal ones.
  \item \textbf{The matroid-rank subclass} (\Cref{sec:mrf}). Write out the
    supermodular-to-matroid reduction properly, in particular the
    parallel-extension step, then attempt basis exchange in the parallel-copy
    matroid. This is where the surrounding literature's machinery is strongest.
  \item \textbf{Build the two tiers first} (Proposition~\ref{prop:twotier}).
    Choose the partition into paid and unpaid agents and only then allocate.
    This is the one natural approach that Theorem~\ref{thm:obstruction} does not
    obviously kill, since it is not incremental in the chores. In the proof of
    Theorem~\ref{thm:binadd} the paid set is exactly the agents holding a
    $\lceil q/n \rceil$-share of the universally disliked chores; identifying
    what plays that role in general is the first question.
  \item \textbf{Climb the milestone ladder.} Binary additive is done
    (Theorem~\ref{thm:binadd}); binary submodular is the next rung, by item~6
    above; general dichotomous is the summit. Note that
    Theorem~\ref{thm:sizeshift} and Remark~\ref{rem:balance} agree on what the
    middle rung needs --- cardinality balancing --- which is what matroid union
    would supply.
  \item \textbf{The mixed model.} Allow each item to be a good or a chore for
    each agent, with all marginals in $\set{-1,0,1}$: the dichotomous
    restriction of the doubly monotone setting of~\cite{KMSTY24}. This is the
    natural common generalisation of~\cite{BKNS22} and of the present paper.
\end{enumerate}

\paragraph{Where this would sit.}
Conjecture~\ref{conj:main} is the fourth corner of the square in
\Cref{sec:intro}: additive goods~\cite{BDNSV20}, dichotomous
goods~\cite{BKNS22} and additive chores~\cite{LMS26} all give one dollar per
agent and $n-1$ in total, and dichotomous chores is what remains. For valuations
that are dichotomous but not additive the best published bound is still that
of~\cite{KMSTY24}, $n-1$ per agent and $n(n-1)/2$ in total from any EF1
allocation, so the conjecture would improve the total by a factor of $n$ ---
the same factor~\cite{BKNS22} gains over~\cite{BDNSV20} on the goods side.

Two cautions follow from the literature check. First,
Theorem~\ref{thm:binadd} is already subsumed by~\cite{LMS26}
(Remark~\ref{rem:binadd-superseded}); what survives of it is the mechanism, not
the statement. Second, \cite{LMS26} appeared in July 2026, which is a reminder
that the surrounding corners are being filled in quickly and that the search
should be repeated before circulation rather than treated as done.
